\section{Introducció i delimitació}

Les aplicacions web executen sovint operacions que no cal completar dins del
cicle immediat d'una petició. La generació d'un document, l'enviament d'una
comunicació, l'actualització d'un índex o el càlcul d'un informe es poden
representar com a tasques en segon pla. Aquest desacoblament evita mantenir
l'usuari esperant, però introdueix un sistema format per productors, cues i
processos de treball que comparteixen una capacitat limitada. Quan s'acumulen
tasques, la política de planificació condiciona qui espera, durant quant temps i
amb quin aprofitament dels recursos.

Aquest estat de l'art estudia els conceptes necessaris per dissenyar una
plataforma web que permeti representar aquest problema amb dades sintètiques.
La revisió se centra en quatre qüestions: les polítiques que decideixen quina
tasca s'executa, les mètriques que permeten comparar-les, el mecanisme
experimental necessari per obtenir resultats reproduïbles i les eines existents
que resolen problemes relacionats. L'objectiu no és demostrar que una política
sigui superior, sinó delimitar com es pot avaluar sota diferents condicions de
càrrega.

La literatura sobre planificació inclou processadors, xarxes, fabricació,
computació d'altes prestacions i infraestructures cloud. Aquests dominis no són
equivalents a una cua de tasques backend. Per això només es recullen resultats
externs quan aporten un concepte transferible, com la inanició, l'equitat,
l'envelliment o el disseny d'experiments. Queden fora els planificadors de temps
real amb garanties estrictes, els workflows amb dependències complexes i la
selecció automàtica de polítiques mitjançant aprenentatge automàtic.

També cal separar les cues operatives dels entorns experimentals. BullMQ o
RabbitMQ executen treball real i incorporen restriccions de persistència,
consum, relliurament i concurrència. En el TFM, en canvi, la simulació serà un
mecanisme intern per comparar polítiques sense esperar el temps real. El
resultat principal serà una plataforma web per definir escenaris, iniciar
comparacions i interpretar-ne les mètriques. El motor donarà suport a aquest
flux, però no substituirà el disseny i el desenvolupament de l'aplicació web.
